\begin{textsample}{POS Dim 9 – global\_south\_2025\_06 – Score 2.00 – global\_south\_2025\_06/126471551.txt}  \label{ex:f9_pos_187}
Stop tolerating terror Muhammad Umar Monday, Jun 30, 2025 Twelve days of missiles, drones, and air strikes left Iranian cities scarred, Israeli towns shaken and the entire region teetering on the edge of a far wider war.
It took direct US intervention to stop the fighting again.
President Trump ' s ceasefire announcement on June 24 brought a fragile pause, but not before Israel had struck deep into Iran and Tehran had responded with missiles that landed near Beersheba.
What came next surprised even longtime Trump watchers.
Standing on the White House lawn, he called Israel out."I ' ve got to get Israel to calm down", he said."I don't like the fact that Israel went out this morning at all."The language was unusually direct.
And when Israeli jets took off for another strike, Trump issued an all-caps warning on Truth Social telling them to abort.
By his account, they did.
For decades, Washington has tiptoed around @ @ @ @ @ @ @ @ @ @ Eastern politics rather than a central problem.
But the events of this war have forced a reckoning.
Iran, crippled by sanctions, battling unrest at home, was willing to stand down once Israel stopped bombing.
It was Israel that kept pushing, even after claiming victory.
The aggression served no strategic value to the US.
It undermined a ceasefire brokered in good faith, put US allies in the crossfire, and dragged the region closer to chaos.
This realisation arrives at a crucial time.
Only weeks earlier, Trump had toured the Gulf, receiving a hero ' s welcome and signing historic economic deals. \textbf{Saudi} Arabia pledged \$142 billion in defence purchases.
Qatar Airways signed a \$96 billion deal with Boeing.
The UAE unveiled a trillion-dollar investment roadmap in data centres and AI ventures in collaboration with US firms.
These are commitments, signalling that the Arab world, long seen through the lens of conflict and crisis, is betting on a different future : one built on innovation, clean energy, digital @ @ @ @ @ @ @ @ @ @ ' s good for the region and great for America.
But there ' s one force pulling in the opposite direction.
One country that continues to test the limits of impunity.
One state whose disregard for restraint is now a liability for peace.
And no, it isn't Iran.
It ' s Israel.
The 12-day war showed the world what many in the region have long known : Israel ' s military dominance has evolved into a doctrine of unchecked force.
More than 430 Iranian civilians were reported killed in the first nine days and over 30 Israelis also died in retaliatory strikes.
Markets shuddered.
Oil prices spiked.
The investment optimism Trump generated in the Gulf suddenly looked premature.
Israel ' s actions put stability, diplomacy and prosperity at risk.
If President Trump is serious about earning a peacemaker title, this is his moment.
Not to distance himself from Israel, but to lead it towards accountability.
Here ' s how : halt all new offensive-grade weapons transfers to Israel until it agrees to verifiable limits @ @ @ @ @ @ @ @ @ @ to Gaza and the West Bank, under international supervision, with clear consequences for obstruction.
Establish a firm deadline for peace talks that include Palestinian representatives and are premised on recognition of a Palestinian state within internationally accepted parameters.
These are not radical demands.
They are the bare minimum to restore the legitimacy of any US-led peace effort.
Pakistan, for its part, should now offer guidance.
As a nation that has navigated complex counterterrorism challenges and championed Muslim causes at global forums, Pakistan is well-positioned to help.
It can provide diplomatic support, verification frameworks and regional insight.
It can help bridge the trust gap between Washington and the broader Muslim world.
That assistance, offered in good faith, can help Trump turn a tactical ceasefire into a strategic breakthrough.
None of this is about obliterating Israel.
That accusation is tired and dishonest.
The issue is not Israel ' s existence.
It is Israel ' s behaviour.
For too long, its occupation, bombardment and blockade of Arab civilians, @ @ @ @ @ @ @ @ @ @ justified as a necessary defence.
But justice and defence are not the same thing.
What millions in the Arab world, and increasingly many Americans, demand is accountability.
Not destruction, accountability.
That starts with ending the impunity that has allowed Israeli terror to continue unchecked since 1948.
Recognition of Palestine is a moral and political necessity for a just and lasting peace.
It is the starting point for a new Middle East where security is shared, not hoarded and where innovation is possible because fear is no longer the \textbf{currency} of diplomacy.
Trump has built his legacy on ' big deals '.
But the biggest deal available to him now is not an arms sale or AI partnership.
It ' s a new framework for peace that recognises Israeli excess, elevates Arab dignity, and finally makes good on the promise of a two-state solution.
He doesn't have to abandon his allies.
He just has to make them act like allies, not arsonists.
If he succeeds, awards will follow @ @ @ @ @ @ @ @ @ @ The writer is a non-resident fellow at the CISS.
He posts/tweets @umarwrites

% matched lemmas: currency, saudi
\end{textsample}
